\documentclass[11pt]{article}

\usepackage[margin=1in]{geometry}
\usepackage{amsmath,amssymb}
\usepackage{booktabs}
\usepackage{tabularx}
\usepackage{array}
\usepackage{enumitem}
\usepackage[T1]{fontenc}
\usepackage{hyperref}

\newcommand{\code}[1]{\texttt{#1}}
\newcommand{\keff}{k_{\mathrm{eff}}}
\newcolumntype{L}{>{\raggedright\arraybackslash}X}
\emergencystretch=3em

\title{Fluctuating hydrodynamics of macroscopic fluctuation theory\\
in \code{src\_deankow/MFT}}
\author{FHDeX}
\date{}

\begin{document}
\maketitle

\noindent This note describes the stochastic PDE solved in
\code{src\_deankow/MFT}, its discretization, the algorithm, and two test
cases: the non-interacting gas and the symmetric simple exclusion process
(SSEP). The spatial and temporal discretization are those of the
Dean--Kawasaki solver in \code{src\_deankow/interaction}, without
particles, mesh refinement or the nonlocal interaction. The code is built
from \code{exec/dean\_kow/MFT}.

\section{Model}

The density $\rho(\mathbf{x},t)$ obeys
\begin{equation}\label{eq:mft}
  \partial_t \rho = \nabla\cdot\Big[ D(\rho)\,\nabla\rho
     + \sqrt{\frac{\sigma(\rho)}{\Lambda}}\;\boldsymbol{\eta} \Big],
\end{equation}
where $\boldsymbol{\eta}$ is space-time white noise with
$\langle \eta_\alpha(\mathbf{x},t)\,\eta_\beta(\mathbf{x}',t')\rangle
 = \delta_{\alpha\beta}\,\delta(\mathbf{x}-\mathbf{x}')\,\delta(t-t')$.
The diffusivity $D(\rho)$ and the mobility $\sigma(\rho)$ are given in the
inputs as \code{amrex::Parser} expressions in the variable \code{rho}
(\code{D\_rho}, \code{sigma\_rho}); $\Lambda$ is the input \code{Lambda}.

\paragraph{Meaning of $\Lambda$.}
For independent Brownian particles of diffusivity $D$, the number density
$n$ obeys the Dean--Kawasaki equation
$\partial_t n = D\nabla^2 n + \nabla\cdot(\sqrt{2Dn}\,\boldsymbol\eta)$.
Writing $n=\Lambda\rho$ gives \eqref{eq:mft} with $D(\rho)=D$ and
$\sigma(\rho)=2D\rho$. So $\Lambda\rho$ is a particle number density and a
cell of volume $\Delta V$ holds $\Lambda\rho\,\Delta V$ particles. For a
lattice gas, $\rho$ is the occupation fraction and $\Lambda\Delta V$ is the
number of lattice sites in the cell. The noise is weak, and the equation is
meaningful as a coarse-grained model, when $\Lambda\rho\,\Delta V\gg1$
(for the SSEP also $\Lambda(1-\rho)\Delta V\gg1$). The relative
fluctuation of a cell density is of order $(\Lambda\rho\,\Delta V)^{-1/2}$.

\paragraph{Examples.}
\begin{center}
\begin{tabular}{lll}
\toprule
model & $D(\rho)$ & $\sigma(\rho)$ \\
\midrule
non-interacting gas (Dean--Kawasaki, $N=\Lambda$) & $D$ & $2D\rho$ \\
SSEP & $1$ & $2\rho(1-\rho)$ \\
\bottomrule
\end{tabular}
\end{center}

\paragraph{Equilibrium fluctuations.}
Linearizing \eqref{eq:mft} about a uniform state $\bar\rho$, the Fourier
modes of $\delta\rho=\rho-\bar\rho$ are Ornstein--Uhlenbeck processes
with relaxation rate $\bar D k^2$, $\bar D = D(\bar\rho)$, and noise strength
$\bar\sigma k^2/\Lambda$, $\bar\sigma=\sigma(\bar\rho)$. The code reports the
static structure factor normalized as $|\hat n|^2/N$ with
$n=\Lambda\rho\,\Delta V$ the cell counts and $N=\Lambda\bar\rho V$:
\begin{equation}\label{eq:S}
  S(k) = \frac{\Lambda}{\bar\rho V}\,\big\langle|\delta\hat\rho(\mathbf{k})|^2\big\rangle
       = \frac{\sigma(\bar\rho)}{2\bar\rho\,D(\bar\rho)},
\end{equation}
where $\delta\hat\rho(\mathbf{k}) = \int_V \delta\rho\, e^{-i\mathbf{k}\cdot\mathbf{x}}\,d\mathbf{x}$.
The intermediate scattering function is
$F(k,\tau) = \frac{\Lambda}{\bar\rho V}\langle\delta\hat\rho(\mathbf{k},t+\tau)\,
\delta\hat\rho(\mathbf{k},t)^*\rangle = S(k)\,e^{-\bar Dk^2\tau}$.
The non-interacting gas has $S=1$ and the SSEP has $S=1-\bar\rho$.

\section{Discretization}

\paragraph{Grid.}
Cell-centred densities $\rho_{\mathbf{i}}$ on a uniform grid with spacing
$\Delta x_d$ and cell volume $\Delta V=\prod_d\Delta x_d$. Fluxes live on the
faces. In one direction (the others are analogous), the face $i-\tfrac12$
separates cells $i-1$ and $i$.

\paragraph{Face densities.}
The deterministic flux uses the arithmetic average
$\rho_f = \tfrac12(\rho_{i-1}+\rho_i)$. The noise amplitude uses the average
of the non-negative parts,
$\rho_f^{+} = \tfrac12\big(\max(\rho_{i-1},0)+\max(\rho_i,0)\big)$, as the
noise in \code{interaction} does, and $\sigma$ is clipped at zero:
\begin{equation}\label{eq:flux}
  F_{i-1/2} = D(\rho_f)\,\frac{\rho_i-\rho_{i-1}}{\Delta x}
     + \sqrt{\max\!\big(\sigma(\rho_f^{+}),0\big)}\;W_{i-1/2}.
\end{equation}
The inner clip turns the noise off next to a cell with negative density for
any $\sigma$ with $\sigma(0)=0$, including models with $\sigma>0$ for
$\rho<0$, such as $\sigma=\rho^2$. The outer clip handles $\sigma<0$ for
non-negative $\rho$, e.g.\ the SSEP for $\rho>1$.

\paragraph{Noise.}
$W_{i-1/2}$ are independent Gaussians with mean $0$ and variance
$1/(\Lambda\,\Delta t\,\Delta V)$, i.e.\ white noise averaged over a cell
volume and a time step, with the factor $1/\Lambda$ included. They are drawn
by \code{MultiFabFillRandom} once per face and step. For $\sigma=2D\rho$ and
$\Lambda=N$ this is the noise of \code{interaction}, whose variance
$2D/(N\,\Delta t\,\Delta V)$ multiplies $\phi$.

\paragraph{Walls.}
Boundaries are periodic (\code{bc = 0}) or zero-density Dirichlet walls
(\code{bc = 3}), treated as in \code{interaction}. On a wall face, the ghost
density is $2\rho_b-\rho$ with $\rho_b=0$; both face densities are set to
$\rho_b$; and the noise is multiplied by $\sqrt2$ (half-cell distance to the
wall).

\paragraph{Time stepping.}
Euler--Maruyama,
\begin{equation}\label{eq:update}
  \rho_{\mathbf i}^{n+1} = \rho_{\mathbf i}^{n}
   + \Delta t \sum_d \frac{F^n_{\mathbf i+\mathbf e_d/2}-F^n_{\mathbf i-\mathbf e_d/2}}{\Delta x_d},
\end{equation}
with all fluxes evaluated from $\rho^n$. The scheme is conservative: with
periodic boundaries $\sum_{\mathbf i}\rho_{\mathbf i}\Delta V$ is constant to
round-off.

\paragraph{Ensemble directions.}
Directions with \code{is\_ensemble\_dir = 1} carry no flux and no noise, so
each row of cells along them is an independent realization. For example,
\code{is\_ensemble\_dir = 0 1} in 2D runs $n_y$ independent 1D systems.

\paragraph{Time step.}
Unless \code{adv.fixed\_dt} is given,
\begin{equation}
  \Delta t = \mathrm{cfl}\,\frac{2}{D_{\max}\,\sum_{d\notin\mathrm{ens}} 4/\Delta x_d^2},
\end{equation}
the forward-Euler limit of the diffusion operator times \code{adv.cfl}.
Here $D_{\max}$ is the maximum of $D$ over the cell values (and $D(0)$ if
there is a wall). $\Delta t$ grows by at most a factor $1.1$ per step.

\section{Algorithm}

Each step of \code{main.cpp}:
\begin{enumerate}[nosep]
  \item swap $\rho^{\mathrm{old}}\leftrightarrow\rho^{\mathrm{new}}$ and fill
        ghost cells (periodic and parallel);
  \item compute $\Delta t$ (\code{est\_time\_step});
  \item draw $W$ on every face of the non-ensemble directions
        (\code{advance\_rho});
  \item compute the face fluxes \eqref{eq:flux} (\code{compute\_flux\_x/y/z}
        in \code{mykernel.H});
  \item update the cells with \eqref{eq:update} (\code{update\_rho});
  \item print $\int\rho$ before and after the step;
  \item every \code{struct\_fact\_int} steps after \code{n\_steps\_skip},
        accumulate $|\mathrm{FFT}(\rho)|^2$ for $S(\mathbf k)$; feed the ISF;
  \item every \code{amr.plot\_int} steps, write the plotfile ($\rho$),
        $S(\mathbf k)$ (\code{plt\_SF*} and the shell average
        \code{sf\_kshell\_<step>.txt}) and the ISF (\code{<isf\_plot\_file>\_1D\_q<q>.dat}).
\end{enumerate}
The initial density is the uniform \code{rho0} or the piecewise-constant
profile along $x$ given by \code{init\_profile\_x/\_val}, used as given (no
renormalization). The normalization \eqref{eq:S} is obtained by passing
\code{StructFact} the scaling $\bar\rho/(\Lambda\Delta V)$, with
$\bar\rho$ the mean initial density; the ISF uses the same factor, so
$F(k,0)=S(k)$. The ISF class is that of \code{interaction}: rolling windows
of x modes, either per x pencil (\code{isf\_do\_1D = 1}, one realization per
pencil) or from the $k_y=k_z=0$ modes of the full FFT.

\section{Discrete equilibrium}

For the linearized scheme on a periodic grid, a Fourier mode obeys
$\delta\hat\rho^{n+1} = (1-a)\,\delta\hat\rho^{n} + \xi^n$, with
$a = \Delta t\,\bar D\,\keff^2$ and
$\keff^2 = \sum_d (4/\Delta x_d^2)\sin^2(k_d\Delta x_d/2)$. The divergence of
the face noise has the same symbol $\keff^2$ as the discrete Laplacian, so
the discrete fluctuation--dissipation balance holds up to the time step
error:
\begin{equation}\label{eq:Sdt}
  S_{\Delta t}(\mathbf k) = \frac{S}{1-a/2},\qquad
  F_{\Delta t}(\mathbf k,\tau) = S_{\Delta t}(\mathbf k)\,(1-a)^{\tau/\Delta t}.
\end{equation}
The Euler--Maruyama bias $1/(1-a/2)$ is largest at the grid scale. For the
time step above, $a\le\mathrm{cfl}$, so at most $1/(1-\mathrm{cfl}/2)$.

\section{Tests}

\subsection{Non-interacting gas and SSEP in 2D}

Input \code{exec/dean\_kow/MFT/inputs\_mft}: a periodic unit square with
$64^2$ cells, $\bar\rho=0.3$, $D=1$, $\Lambda=10^6$ ($\Lambda\bar\rho\Delta V\approx73$
particles per cell), $\mathrm{cfl}=0.1$ ($\Delta t=6.1\times10^{-6}$),
20000 steps, $S$ sampled every 10 steps after step 8000 (1200 samples).
The non-interacting gas uses \code{sigma\_rho = "2.*rho"}, the SSEP
\code{sigma\_rho = "2.*rho*(1.-rho)"}.

\begin{center}
\begin{tabularx}{\linewidth}{lccL}
\toprule
 & theory $S$ & shells $20<k<120$ & $\Lambda\,\mathrm{var}(\rho)\,\Delta V/\bar\rho$ (all $k$) \\
\midrule
non-interacting & $1$ & $1.025$ & $1.071$ \\
SSEP & $1-\bar\rho=0.7$ & $0.712$ & $0.736$ \\
\bottomrule
\end{tabularx}
\end{center}

\noindent The excess over the theory is the Euler--Maruyama bias
\eqref{eq:Sdt}: its average over all modes is $\approx1.05$ here, which
accounts for the real-space variance, and the highest shell
($k\approx195$) gives $1.063$ and $0.743$ against $1.053$ and $0.737$
predicted. The lowest shells hold few modes and relax on
$1/(\bar D k^2)\approx0.025$, comparable to the sampling window, so they
are noisier. With periodic boundaries the total mass changes only at
round-off ($\sim10^{-15}$).

\subsection{1D ensemble: intermediate scattering function of the SSEP}

Input \code{inputs\_mft\_1d\_ensemble\_isf}: $1000$ independent periodic 1D
systems (\code{is\_ensemble\_dir = 0 1}) of length $L=256$ with
$\Delta x=1$, SSEP at $\bar\rho=0.5$, $\Lambda=100$,
$\Delta t = 0.05$ ($\bar D\Delta t/\Delta x^2=0.05$), 30000 steps. The ISF
of the modes $q=4,12,20,28$ ($k=2\pi q/L$) is averaged over pencils and
over 191 time windows, starting at step 10000 with lags up to $\tau=50$.

\begin{center}
\begin{tabular}{rccccc}
\toprule
$q$ & $F(k,0)$ & $F_{\Delta t}(k,0)$ & $F(k,5)$ & $F_{\Delta t}(k,5)$ & $\max_\tau|F-F_{\Delta t}|$ \\
\midrule
 4 & 0.496 & 0.500 & 0.473 & 0.477 & 0.004 \\
12 & 0.498 & 0.501 & 0.323 & 0.325 & 0.004 \\
20 & 0.501 & 0.503 & 0.152 & 0.153 & 0.003 \\
28 & 0.502 & 0.506 & 0.050 & 0.051 & 0.003 \\
\bottomrule
\end{tabular}
\end{center}

\noindent Measured values agree with \eqref{eq:Sdt} to within about 1\% of
$S$ at all lags, and $F(k,0)$ agrees with the static $S=1-\bar\rho=0.5$.

\end{document}
